\documentclass[a4paper]{article}
\usepackage{color,epsfig}
\usepackage{makecell}
\newcolumntype{C}[1]{>{\centering\arraybackslash}m{#1}}
\usepackage{array}
\usepackage{longtable}
\usepackage{dirtytalk}
\usepackage{afterpage}
\usepackage{xcolor}
\newcommand{\pdif}[2]{\frac{\partial #1}{\partial #2}}
\newcommand{\oline}[1]{ \overline{#1}}
\definecolor{RED_FLAG}{rgb}{0.8,0,0}
\definecolor{BLUE_FLAG}{rgb}{0,0.05,0.99}
\definecolor{Olivegreen}{rgb}{0,.4,0}
\begin{document}

{\Large \bf {Response to editor comments}}
\vspace{1cm}

{\textcolor{Olivegreen}{
We appreciate the comments from the editor and two reviewers. We have carefully examined all the concerns, which are explained in more detail in two response letters.}}
%%%%%%%%%%%%%%%%%%%%%%%%%%%%%%%%%%%%%%%%%%%%%%
\begin{itemize}
	\item\textit{Your manuscript ID 230570091 entitled "Reynolds number effect on the Lagrangian evolution of hairpin vortices in turbulent channel flow" which you submitted to the Journal of Turbulence, has been reviewed and I include the comments of the reviewer(s) at the bottom of this letter. You will see that they have raised significant issues which must be addressed before we can consider it further. I do hope that you would be willing to undertake the additional work and resubmit the manuscript.}
\vspace{.2cm}

Thank you for your valuable feedback and guidance regarding the review reports. We have addressed all the significant issues raised by the reviewers. In addition, we have gone through all the minor comments and implemented them in the revised manuscript.

%%%%%%%%%%%%%%%%%%%%%%%%%%%%%%%%%%%%%%%%%%%%%%
\item \textit{In particular, I would like to draw your attention to the comments of Reviewer 1.   Indeed, two hairpins are insufficient to draw any conclusion, and it is, in my opinion, necessary to show some generality of the results shown.}
\vspace{.2cm}

Thank you for your kind suggestion. We have identified five hairpin vortices for each Reynolds number to increase the sample size which enhances the generality of the discussions. The details of these new structures are given in the response letter to reviewer 1.



%%%%%%%%%%%%%%%%%%%%%%%%%%%%%%%%%%%%%%%%%%%%%%
\item \textit{The second point the reviewer makes is also valid.  A better description of the extraction of the core vortex line must be provided (and if it is indeed identified visually the utility of the proposed method is greatly diminished).}
\vspace{.2cm}

Thank you for your kind suggestion. We have provided a more comprehensive description of the extraction of the core vortex line. The procedure for determining the core vortex line is not visual, only the initial point of the extraction method is selected manually, then the rest of the process is completely automatic. A complete description of the procedure is provided in response letters to both reviewers.

%%%%%%%%%%%%%%%%%%%%%%%%%%%%%%%%%%%%%%%%%%%%%%
\item \textit{The modifications requested by Reviewer 2 are more straightforward, and I am sure you can change your manuscript accordingly.}
\vspace{.2cm}

We thank you for your positive feedback. We have changed the revised manuscript according to the modifications requested by Reviewer 2.

%.................................................................
\end{itemize}
%.................................................................

%\bibliographystyle{tfnlm}
%\bibliography{Response}
%.................................................................
\end{document}


